% ═══ Related Work ═══
\section{Related Work}
\label{sec:related}

\subsection{Coevolutionary contact prediction}
The inference of residue--residue contacts from multiple sequence alignments
(MSAs) has evolved through several generations of methodology.  Early
approaches based on mutual information~\cite{dunn2008} suffered from
phylogenetic and entropic biases that produced high false-positive rates.
Direct Coupling Analysis (DCA) addressed this by modelling the joint sequence
distribution through a Potts Hamiltonian whose coupling parameters capture
direct (rather than transitive) correlations~\cite{weigt2009}.  The mean-field
approximation of Morcos et al.\cite{morcos2011} made DCA computationally
tractable for large alignments by replacing the full partition function with a
covariance-based estimate.  Sparse inverse covariance estimation, as
implemented in PSICOV~\cite{jones2012}, further improved precision by
promoting sparsity in the precision matrix through $\ell_1$ regularisation.
Pseudo-likelihood maximisation~\cite{ekeberg2013} and fast belief propagation
implementations such as CCMpred~\cite{seemayer2014} subsequently refined the
inference procedure.  GREMLIN~\cite{kamisetty2013} combined DCA with
protein-structure priors.  EVmutation~\cite{hopf2017} extended the framework
to predict mutation effects on fitness rather than contacts alone.

More recently, end-to-end deep learning approaches such as
AlphaFold2~\cite{jumper2021} have integrated coevolutionary signal with
structural priors inside attention-based architectures, achieving near-experimental
accuracy for many targets.  However, these models produce dense internal
representations that resist decomposition into interpretable components.

\subsection{Formal methods in computational biology}
The application of interactive theorem provers to computational biology
remains limited.  Existing formalisations focus on algorithm correctness for
phylogenetic inference~\cite{lean4}, genome assembly graph properties, and
sequence alignment dynamic programming recurrences.  No prior work has
established machine-checked proofs that the encoding layer connecting amino-acid
alphabets to numerical representations preserves the structural properties
required for coevolutionary analysis.

\subsection{Karnaugh maps beyond digital logic}
Karnaugh maps were introduced as a graphical method for minimising Boolean
functions~\cite{karnaugh1953}.  Their application outside digital circuit
design has been explored in synthetic biology, where Boolean gates model gene
regulatory networks~\cite{marchisio2011}, and in the analysis of genetic code
symmetries~\cite{petoukhov2025}.  Our work extends this tradition by using
Karnaugh maps not merely as a visual representation but as an analytical
engine: the Quine--McCluskey minimiser extracts testable coevolutionary rules
from thresholded frequency data, and the resulting circuits are validated
against native protein structures.
